% The Clean Two-Query Threshold for Unitary Synthesis -- LaTeX source.
% Content matches paper/paper.md version 1.0.4 (the markdown is the GitHub-rendered copy).
% Compile:  pdflatex paper.tex, three passes (or latexmk -pdf paper.tex)
\documentclass[11pt,a4paper]{article}

\usepackage[utf8]{inputenc}
\usepackage[T1]{fontenc}
\usepackage{lmodern}
\usepackage{amsmath,amssymb,amsthm}
\usepackage[margin=1.1in]{geometry}
\usepackage{booktabs}
\usepackage{enumitem}
\usepackage{microtype}
\usepackage{xcolor}
\usepackage[colorlinks=true,linkcolor=blue!50!black,citecolor=blue!50!black,urlcolor=blue!50!black]{hyperref}
\hypersetup{pdftitle={The Clean Two-Query Threshold for Unitary Synthesis},
pdfauthor={Seth Douglas},pdfsubject={Clean two-query unitary synthesis},
pdfkeywords={quantum query complexity, unitary synthesis, Gaussian mean width}}

\theoremstyle{plain}
\newtheorem{theorem}{Theorem}[section]
\newtheorem{proposition}[theorem]{Proposition}
\newtheorem{corollary}[theorem]{Corollary}
\newtheorem{lemma}[theorem]{Lemma}
\theoremstyle{definition}
\newtheorem{definition}[theorem]{Definition}
\theoremstyle{remark}
\newtheorem{remark}[theorem]{Remark}

\DeclareMathOperator{\Tr}{Tr}
\DeclareMathOperator{\ran}{ran}
\DeclareMathOperator{\Pack}{Pack}
\DeclareMathOperator{\Ham}{Ham}
\DeclareMathOperator{\diag}{diag}
\DeclareMathOperator{\rank}{rank}
\DeclareMathOperator{\Gr}{Gr}
\DeclareMathOperator{\polylog}{polylog}
\DeclareMathOperator{\poly}{poly}
\DeclareMathOperator{\Ad}{Ad}
\DeclareMathOperator{\op}{op}
\newcommand{\T}{\mathbb T}
\newcommand{\C}{\mathbb C}
\newcommand{\R}{\mathbb R}
\newcommand{\E}{\mathbb E}
\newcommand{\Fro}[1]{\|#1\|_F}
\newcommand{\norm}[1]{\left\|#1\right\|}
\newcommand{\Lawful}{\mathcal L}

% a plain, package-free call-out box
\newenvironment{callout}[1]{\par\medskip\noindent\begin{quote}\small\textbf{#1}\ }{\end{quote}\medskip}

\title{The Clean Two-Query Threshold for Unitary Synthesis}
\author{Seth Douglas}
\date{Version 1.0.4 --- 2026-09-13}

\begin{document}
\maketitle

\begin{abstract}
Aaronson and Kuperberg asked whether every $n$-qubit unitary can be implemented by a
polynomial-size quantum circuit making few queries to a classical oracle. At one query the
answer is known to be no. We study two sequential queries, allowing arbitrary
workspace but requiring its approximate return to a fixed state.

We settle the two-query case in the clean model, where the workspace must be returned to
its initial state, by pinning both of its resource thresholds up to a logarithm in the
error. Write $N=2^n$.
\begin{itemize}[nosep]
\item \textbf{Width.} For any fixed two-query architecture $T_{x,y}=B\,D_x\,W\,D_y\,C$ with
arbitrary contractions $B,W,C$ and independent unimodular phase selectors $x,y$, the
Gaussian mean width of the realizable family is at most $2\min(N,Q)\sqrt{\log 2Q}$, where
$Q$ is the dimension of the register the oracle acts on, ancillas included. The supremum
over selectors is absorbed pointwise; no chaining is used.
\item \textbf{Lower bounds.} Consequently a universal clean two-query circuit needs
$\Omega(2^n)$ qubits, and, by an address-compression lemma, an oracle input length of
$\Omega(2^n)$ bits regardless of workspace.
\item \textbf{Matching upper bound.} A permutation-transport construction implements every
unitary cleanly with two queries, $O(2^n\log(1/\varepsilon))$ qubits and
$O(2^n\log(1/\varepsilon))$ address bits, at error $\varepsilon$.
\end{itemize}
So the clean two-query qubit threshold and address-length threshold are both $\Theta(2^n)$
up to the factor $\log(1/\varepsilon)$. We also show that the width bound is tight at register
dimension $Q=\Theta(N^2)$, explain why the one-query argument does not extend, and prove a transfer
theorem reducing the garbage-model two-query problem to a clean four-insertion width bound,
which remains open. This is a sufficient route, not a proof that every garbage-model
method must use four insertions.
\end{abstract}

\begin{callout}{Verification note.}
Selected identities and finite instances are exercised by scripts in
\texttt{verify/} of the \href{https://github.com/Apsiape/two-query-unitary-synthesis}{accompanying repository}, and
\texttt{STATUS.md} there records, claim by claim, which checks exist, what has been
independently audited and what has not. These checks do not certify the complete
analytic proof. File names in this paper are paths in that
repository.
\end{callout}

\section{Introduction}

\subsection{The problem}

The unitary synthesis problem of Aaronson and Kuperberg \cite{AK07} asks whether every
$n$-qubit unitary can be implemented by a polynomial-time quantum algorithm with few queries
to a classical oracle. The known landmarks are:
\begin{itemize}[nosep]
\item \textbf{\cite[Thm~6.7]{AK07}} one query, exact model: at most $4^N$ distinct
unitaries, at most $2N$ relevant Boolean table entries (not address bits);
\item \textbf{\cite{LMW24}} one query, or polynomially many parallel (non-adaptive) queries,
in the approximate/discard model, oracle input length $o(2^n)$;
\item \textbf{\cite{DLM26}} explicit one-query separations, and Question~1.4 posing the
general-depth problem;
\item \textbf{\cite{BY26}} an oracle input-length bound $(2-o(1))\log d$ valid at any query
count;
\item \textbf{\cite{Ros26}} an upper bound: $O(2^{n/2})$ clean queries suffice, with
$\poly(n)$ ancillas.
\end{itemize}
This paper proves a workspace-uniform capacity bound for two sequential queries in the
clean model and a matching fixed-error resource threshold. Related restricted work
includes Banerjee's exact every-oracle-clean counting analysis \cite{Ban25} and Huang's
1.5-query study \cite{Hua25}. Those models and statements should not be identified with
the present approximation-and-address theorem; none of our proofs uses them.

\subsection{The model and the scope of the clean restriction}

A fixed circuit with two oracle calls realizes a target $U$ if some oracle makes it
implement $U$ on the target register and, in the \textbf{clean} model, return every
auxiliary register to a fixed state. Nothing is required of the circuit for other oracles;
in particular, unlike the setting of \cite[Thm~6.7]{AK07}, whose footnote~13 notes that the
theorem needs the circuit to implement some unitary for every oracle, we impose no such
requirement, and our lower bounds hold without it. The Bernstein--Vazirani one-query universal
synthesizer, a lookup attributed to Yuen in \cite{LMW24}, is not
clean: its workspace ends holding the description string. So the class we bound is strictly
smaller than the discard model of \cite{LMW24}, and \textbf{our conclusions are
correspondingly stronger on a correspondingly smaller class; they neither improve nor
contradict \cite{LMW24} or \cite{BY26}.} Section~\ref{sec:garbage} prices the distance
between the two models.

\subsection{What is new}

\cite[Def.~B.1]{LMW24} already name the quantity we bound: an oracle circuit is $S$-small if
it realizes at most $S$ distinct unitaries, and their Appendix~B proves that small capacity
implies a lower bound. They then observe that counting arguments ``become ineffective once
the adversary has even a small number of ancilla qubits.'' The implication is theirs; the
missing premise is that a two-query architecture's capacity \emph{is} small. This paper
supplies that premise, in a metric-entropy form that degrades only as the square root of
the total qubit count: an adversary needs $\Omega(2^n)$ ancilla qubits to neutralize it.

Two rounds, not two queries. A phase selector $D_x$ is an arbitrary diagonal on the whole
register, and a product of phase oracles applied in one round is again such a diagonal. One
selector factor therefore subsumes arbitrarily many parallel queries, so
Theorem~\ref{thm:width} constrains \emph{two rounds of adaptivity with arbitrarily many
parallel queries per round}.

The obstruction \cite[\S2.5]{LMW24} name for going past one query, bounding the operator
norm of $2^M$ matrices simultaneously, is met by absorbing the entire selector supremum
\emph{pointwise} into a single operator norm (Section~\ref{sec:width}).

Beyond the width bound, this version adds: the extension from real sign words to
independent unimodular selectors, which covers diagonal-unitary oracles and two different
tables; an address-length lower bound that holds with any amount of workspace; the matching
upper bound; and the garbage-model transfer.

\paragraph{Indexing.} \cite[Question~1.4]{DLM26} ask whether synthesizing $F_tH\cdots F_1$
requires $t$ sequential queries; the case $t$ is a lower bound against $t-1$ queries, so
their proved case $t=2$ is a one-query bound and a two-query bound has the strength of their
case $t=3$. Our result has that strength for Haar-scale targets; we say nothing about their
explicit family.

\section{Setting}\label{sec:setting}

\begin{definition}[two-query architecture]\label{def:arch}
A triple $(B,W,C)$ with $C$ a $Q\times N$ contraction, $W$ a $Q\times Q$ contraction and $B$
an $N\times Q$ contraction, realizing
\[
T_{x,y}\;=\;B\,D_x\,W\,D_y\,C,\qquad D_x=\diag(x),\ D_y=\diag(y),\qquad x,y\in\T^Q,
\]
where $\T=\{z\in\C:|z|=1\}$. The \textbf{repeated real sign word} is the special case
$x=y=g\in\{\pm1\}^Q$, written $T_g$; it is the shape of a two-query circuit with one $\pm1$
phase oracle. Independent $x,y$ cover diagonal-unitary oracles and circuits that query two
different tables.
\end{definition}

Ancillas are inside $Q$: an architecture on $n$ target qubits and $a$ ancillas has
$Q=2^{n+a}$, and we write $w=\log_2Q$ for the total qubit count. The bounds below are stated
in $\log(2Q)$, i.e.\ in $w$; Theorem~\ref{thm:address} then separates the oracle input length
from the workspace.

\paragraph{From oracle circuits to Definition~\ref{def:arch}.}
Let $\iota_{\rm in},\iota_{\rm out}:\C^N\to\C^Q$ append the fixed initial
and required final auxiliary states. A two-query circuit has full unitary
$V_f=A_2O_fA_1O_fA_0$, with target-independent unitary gates $A_j$.
For the usual Boolean query $O_f|a,b,r\rangle=|a,b\oplus f(a),r\rangle$,
conjugating its answer qubit by a Hadamard gives
$D_f|a,b,r\rangle=(-1)^{bf(a)}|a,b,r\rangle$; the spectator register $r$
is unchanged. Absorb these fixed Hadamards into the adjacent $A_j$, so that
$V_f=A_2D_fA_1D_fA_0$. Its clean-output block is
\[
 T_f=\iota_{\rm out}^*V_f\iota_{\rm in}
     =B D_f W D_f C,\qquad
 B=\iota_{\rm out}^*A_2,\quad W=A_1,\quad C=A_0\iota_{\rm in}.
\]
These are fixed contractions of the stated dimensions. The physical diagonal
entries have repetitions and a fixed $+1$ sector; allowing arbitrary independent
$x,y\in\T^Q$ only enlarges the family. If the successful oracle satisfies the
full-output approximate-clean guarantee
$\|V_f\iota_{\rm in}-\iota_{\rm out}U\|_{\op}\le\gamma$, then
\[
 \|T_f-U\|_{\op}
 =\|\iota_{\rm out}^*(V_f\iota_{\rm in}-\iota_{\rm out}U)\|_{\op}
 \le\gamma.
\]
This is uniform over reference-entangled inputs. It assumes approximate cleanup
of the full output, not merely closeness of the reduced channel after discarding
auxiliaries; the latter is treated separately in Section~\ref{sec:garbage}.

\paragraph{Notation.} $b_p=Be_p\in\C^N$ is the $p$-th column of $B$ and $c_q=e_q^{\top}C\in\C^N$
the $q$-th row of $C$, with \textbf{leverages} $\lambda_p=\|b_p\|^2$, $\mu_q=\|c_q\|^2$,
$\Lambda=\diag(\lambda)$, $M_\mu=\diag(\mu)$, and $P_\lambda,P_\mu$ the projections onto their
supports. Since $B$ and $C$ are contractions, $\sum_p\lambda_p=\Fro{B}^2\le\min(N,Q)$ and
$\sum_q\mu_q=\Fro{C}^2\le\min(N,Q)$.

\begin{proposition}[unitarity criterion]\label{prop:unitarity}
Write $\Psi(x,y)=D_xWD_yC$ and $\Pi=I_Q-B^{*}B$. Then
\[
T_{x,y}^{*}T_{x,y}\;=\;I_N\;-\;(I_N-C^{*}C)\;-\;C^{*}D_y^{*}\,(I_Q-W^{*}W)\,D_yC\;-\;\Psi(x,y)^{*}\,\Pi\,\Psi(x,y).
\]
All three subtracted terms are positive semidefinite, so $T_{x,y}$ is unitary exactly when
$C$ is an isometry, $(I_Q-W^{*}W)D_yC=0$ and $\Pi\Psi(x,y)=0$. When $C$ and $W$ are
isometries the criterion reduces to $\Pi\Psi=0$.
\end{proposition}

\begin{proof}
$T^{*}T=\Psi^{*}B^{*}B\Psi=\Psi^{*}\Psi-\Psi^{*}\Pi\Psi$ and
$\Psi^{*}\Psi=C^{*}D_y^{*}W^{*}WD_yC=C^{*}C-C^{*}D_y^{*}(I-W^{*}W)D_yC$, while
$C^{*}C=I_N-(I_N-C^{*}C)$. Each subtracted term is positive semidefinite because $I-C^{*}C$,
$I-W^{*}W$ and $\Pi$ are, for contractions; and a sum of positive semidefinite operators
vanishes only when each does.
\end{proof}

\begin{callout}{Erratum.}
An earlier draft stated the criterion with only the last term, which is correct only for
isometric $C$ and $W$; with $C=0$, a legal contraction, the shorter criterion is satisfied
while $T=0$. \texttt{verify\_complex\_phase.py}~(e) checks the three-term identity for
contractive $C$ and $W$ and exhibits the failure of the shorter forms.
\end{callout}

\begin{definition}[capacity]\label{def:capacity}
The capacity at scale $\rho$ is $\log\Pack_{\rho\sqrt N}\{T_{x,y}\}$, the logarithm of the
largest family of realizable operators that is pairwise $\rho\sqrt N$-separated in Frobenius
norm. This is the physically meaningful scale for unitary targets, since $\Fro{U}=\sqrt N$.
\textbf{The capacity claim is asserted only in this metric}; a raw, un-normalized reading is
false by a factor $N$ (Corollary~\ref{cor:capacity}).
\end{definition}

\section{Why the one-query argument does not extend}\label{sec:onequery}

\cite[Thm~6.7]{AK07} writes $U(X_{ij})=I+E_i+E_j$ and applies \cite[Lem.~6.6]{AK07} to the
anticommutation relations $E_iE_j^{\dagger}+E_jE_i^{\dagger}=0$. This is pure affineness and
is unavailable at degree two. One might still hope that a combinatorial invariant of the
\textbf{lawful set}, the set of selector words for which the circuit's output is unitary,
controls capacity. Two facts speak against building on that hope, one at depth two and one
beyond it.

\begin{proposition}[at depth two the final layer is invisible]\label{prop:invisible}
For a two-query architecture, write $\Psi(x,y)=D_xWD_yC$ and let $\Lawful$ be its lawful set.
Then the number of distinct realizable targets $\{T_{x,y}:(x,y)\in\Lawful\}$ equals the number
of distinct $\Psi(x,y)$, $(x,y)\in\Lawful$. Given $\Psi$ and $\Lawful$, the count does not
depend on $B$.
\end{proposition}

\begin{proof}
For lawful $(x,y)$, $T=B\Psi$ is unitary, so by Proposition~\ref{prop:unitarity} $\Psi$ is an
isometry and $B$ is isometric on $\ran\Psi(x,y)$. For a contraction $B$ the set of vectors on
which it is isometric is the subspace $\ker(I-B^{*}B)$, so $B$ is injective on the sum of
these ranges over the lawful set, and $B\Psi(x,y)=B\Psi(x',y')$ forces
$\Psi(x,y)=\Psi(x',y')$.
\end{proof}

\begin{proposition}[beyond depth two, pre-final data does not determine the count]\label{prop:pair}
In the three-insertion class $\Psi_g=D_gW_2D_gW_1D_gC$, $T_g=B\Psi_g$, with $Q=8$ and a
one-dimensional target ($N=1$), there are two architectures sharing $C$, $W_1$ and $B$ and
differing only in $W_2$, with identical lawful sets (a linear code of $16$ sign words closed
under pointwise product) and identical pre-final data (the frames $C,\,D_gC,\,D_gW_1D_gC$ on
the lawful set and their Gram kernels), whose numbers of distinct realizable targets are $2$
and $16$.
\end{proposition}

An explicit rational construction and its exhaustive exact certificate are given in
Appendix~\ref{app:pair}. The numerical script \texttt{verify\_decisive\_pair.py}
provides supplementary checks, including Proposition~\ref{prop:invisible}.
Since $N=1$, this example distinguishes scalar matrices, not quantum channels.

\begin{remark}[what this motivates, and nothing more]\label{rem:motivates}
Neither fact is a bound. The first says that at depth two capacity is a property of the
pre-final map restricted to the lawful set, an object whose Boolean structure we have no
tool to count. The second shows that frames ending before the last variable mixer
need not determine the final matrices; it does not contradict injectivity of a fixed
final contraction on its lawful ranges. Both point
away from combinatorial invariants of the lawful set and toward a quantity that sees the
operators $T_{x,y}$ themselves: their metric entropy, which Section~\ref{sec:width} controls
through a Gaussian mean width. A $2$-versus-$16$ gap at $Q=8$ establishes nothing
asymptotically.
\end{remark}

\section{The width theorem}\label{sec:width}

\begin{callout}{Terminology.}
``Width'' means the Gaussian mean width of a set of matrices, in the sense of convex
geometry. \cite[Lem.~4.11]{LMW24} use ``width'' for an unrelated leverage-normalized overlap
statistic.
\end{callout}

Let $G$ be an $N\times N$ matrix with independent standard complex Gaussian entries,
$\E|G_{ij}|^2=1$.

\begin{theorem}[two-query Gaussian mean width, unimodular selectors]\label{thm:width}
For every two-query architecture $(B,W,C)$ of Definition~\ref{def:arch},
\[
w_2\;:=\;\E_G\,\sup_{x,y\in\T^Q}\bigl|\Tr(G^{*}T_{x,y})\bigr|
\;\le\;\Fro{B}\,\Fro{C}\cdot2\sqrt{\log(2Q)}
\;\le\;2\min(N,Q)\sqrt{\log(2Q)}.
\]
The supremum sits inside the expectation; no exchange of supremum and expectation occurs.
\end{theorem}

\begin{proof}
If either endpoint is zero, the assertion is immediate. Otherwise both leverage
supports are nonempty. The following six steps keep the selector supremum inside
the expectation; the absorption inequality is pointwise.
\begin{enumerate}[leftmargin=*]
\item \textbf{Exact expansion.} Since $D_x=\sum_px_pe_pe_p^{\top}$ and likewise for $D_y$,
\[
\Tr(G^{*}T_{x,y})\;=\;\sum_{p,q}x_p\,y_q\,W_{pq}\,\bigl(c_q\,G^{*}\,b_p\bigr)\;=\;x^{\top}M\,y,
\qquad M_{pq}:=W_{pq}\,(c_qG^{*}b_p).
\]
\textbf{The matrix $M$ does not depend on $x$ or $y$.} If $\lambda_p=0$ then $b_p=0$ and row
$p$ of $M$ vanishes; if $\mu_q=0$ then column $q$ vanishes. Restricting to the supports is
therefore exact.

\item \textbf{Leverage normalization.} Put $K:=\Lambda^{-1/2}M\,M_\mu^{-1/2}$ on the supports
and $u=\Lambda^{1/2}x$, $v=M_\mu^{1/2}y$, so that $x^{\top}My=u^{\top}Kv$.

\item \textbf{Pointwise absorption.} For unimodular selectors the norms
$\|u\|^2=\sum_p\lambda_p|x_p|^2=\Fro{B}^2$ and $\|v\|^2=\Fro{C}^2$ are \textbf{constant over
the torus}. Hence, pointwise in $G$,
\[
\sup_{x,y\in\T^Q}\bigl|x^{\top}My\bigr|\;\le\;\Fro{B}\,\Fro{C}\;\|K\|_{\op}.
\]
The supremum over the entire selector torus has been absorbed into one operator norm: no
net, no chaining, no union bound. This is the heart of the proof.

\begin{callout}{Provenance.}
The inequality of this step is not new. It is the diagonal-scaling relaxation of quadratic
maximization over the cube and the torus that runs from Delorme--Poljak and Poljak--Rendl
through Nesterov, Alon--Naor and Charikar--Wirth \cite{DP93,PR95,Nes98,AN06,CW04}; the
degree-$t$ structure is the polynomial method \cite{BBCMW98}, of which \cite[Thm~6.7]{AK07}
is the $t=1$ case. What is ours is the composition: polynomial method, then this relaxation,
then matrix concentration, then Sudakov, run against a synthesis architecture, and the
variance computation of the next step that makes the composition close.
\end{callout}

\item \textbf{$K$ is a matrix Gaussian series.} Writing $G=\sum_{ij}G_{ij}e_ie_j^{\top}$,
\[
K=\sum_{i,j}\overline{G_{ij}}\,A_{ij},\qquad
(A_{ij})_{pq}=\frac{W_{pq}\,(b_p)_i\,(c_q)_j}{\sqrt{\lambda_p\mu_q}}.
\]
Its two variance operators have closed forms. With $\Gamma_B:=\Lambda^{-1/2}B^{*}B\Lambda^{-1/2}$
and $\Gamma_C:=M_\mu^{-1/2}CC^{*}M_\mu^{-1/2}$, both positive semidefinite with unit diagonal on
their supports,
\[
\sum_{ij}A_{ij}A_{ij}^{*}=\overline{\Gamma_B}\circ\bigl(WP_\mu W^{*}\bigr),\qquad
\sum_{ij}A_{ij}^{*}A_{ij}=\overline{\Gamma_C}\circ\bigl(W^{*}P_\lambda W\bigr),
\]
where $\circ$ is the entrywise product. Both closed forms are checked against Monte Carlo in
\texttt{verify\_complex\_phase.py}~(c).

\item \textbf{Both variance operators are dominated by the identity.} If $\Gamma\succeq0$ has
$\Gamma_{pp}\le1$ then $\|\Gamma\circ X\|_{\op}\le\|X\|_{\op}$ for every $X$: writing
$\Gamma_{pp'}=u_p^{*}u_{p'}$ with $\|u_p\|\le1$, the matrix $\Gamma\circ X$ is the compression
of $I\otimes X$ by the contraction $e_p\mapsto u_p\otimes e_p$; $\overline{\Gamma}$ is positive
semidefinite with the same diagonal, so the lemma applies to it as well. Since $WP_\mu W^{*}$
and $W^{*}P_\lambda W$ are positive semidefinite of norm at most one, both variance operators
are positive semidefinite of norm at most one, and they are supported on $P_\lambda$ and
$P_\mu$ respectively. Only $\|W\|\le1$ is used; the middle layer is otherwise arbitrary.

\item \textbf{Matrix Khintchine.} For a real Gaussian series $Z=\sum_k\gamma_kA_k$ with
$d_1\times d_2$ coefficients, Tropp's bound gives $\E\|Z\|\le\sqrt{2v\log(d_1+d_2)}$ with
$v=\max\{\|\sum_kA_kA_k^{*}\|,\|\sum_kA_k^{*}A_k\|\}$ \cite[Thm~4.1.1]{Tro15}. Splitting each
complex Gaussian into real and imaginary parts of variance $1/2$ gives two real series with
variance parameter $v/2\le1/2$, hence $\E\|K\|\le2\sqrt{\log(2Q)}$ by the triangle
inequality. (Treating the two parts as one real series of $2N^2$ terms gives the sharper
$\sqrt2$, which is what the verifier tests; the constant $2$ is all that is needed.) Assemble
with step~3.
\end{enumerate}
\end{proof}

\begin{corollary}[real sign words; the certified form]\label{cor:real}
For the repeated real sign word,
\[
\E_G\sup_{g\in\{\pm1\}^Q}\operatorname{Re}\Tr(G^{*}T_g)\;\le\;2\min(N,Q)\sqrt{\log(2Q)}.
\]
\end{corollary}

\begin{proof}
Take $x=y=g$ and use $|\operatorname{Re}z|\le|z|$.
\end{proof}

The form of this corollary that was independently audited (see \texttt{STATUS.md}) carries
the constant $8$ and is proved by a symmetric weighting: with $d_p=\lambda_p+\mu_p$ and
$D=\diag(d)$ one writes $\operatorname{Re}(g^{\top}Mg)=v^{*}Kv$ with $v=D^{1/2}g$,
$\|v\|^2=\Tr D\le2\min(N,Q)$, and bounds $\E\|K\|\le4\sqrt{\log2Q}$. The asymmetric route
above was also verified in that audit, as the product-form strengthening with
$\Fro{B}\Fro{C}$ in place of $\Tr D$; both routes are valid and incomparable instance by
instance. The deterministic inequalities of the symmetric proof
are formalized in the \texttt{lean/} directory of the repository.

\begin{remark}[the weighting is sufficient, not unique]\label{rem:weighting}
Under $d\mapsto td$ the variance operators scale by $t^{-2}$, so any $d'\succeq\lambda+\mu$
works subject to the trace budget. Balance is necessary: over random contractive
architectures, $d=\lambda$ alone or $d=\mu$ alone fails the domination on a substantial
fraction of instances, while the balanced or the asymmetric normalization succeeds on all
of them.
\end{remark}

\begin{remark}[zero leverage is lossless]\label{rem:zero}
$\lambda_p=0$ kills row $p$ of $M$ and $\mu_q=0$ kills column $q$, so the supports may differ
and $K$ is rectangular in general. Nothing downstream changes; the concentration step stays
at dimension at most $2Q$. The same holds when the middle layer is itself rectangular,
$Q_1\times Q_2$, with a selector of length $Q_1$ on the left and $Q_2$ on the right: nothing in
steps 1 to 6 uses $Q_1=Q_2$, the endpoint norms are as before, and the Khintchine dimension
becomes $Q_1+Q_2$.
\end{remark}

\begin{remark}[provenance of the logarithm]\label{rem:log}
$\sqrt{\log(2Q)}$ is the standard Khintchine logarithm, $2Q$ being Tropp's $d_1+d_2$. It is
the same factor appearing as $\sqrt{2\ln2M}$ in \cite[Lem.~4.11]{LMW24}. Inherited, not
discovered.
\end{remark}

\begin{remark}[the selector need not be real]\label{rem:complex}
An earlier draft proved only the real case and disclaimed any extension to diagonal-unitary
oracles, because its symmetric route rewrote $\operatorname{Re}(g^{\top}Mg)$ as a Hermitian
form, which uses $g\in\R^Q$. The asymmetric route never does this: step~3 uses only that
$|x_p|=|y_q|=1$. The real restriction is gone, and with it the restriction to a single table.
\end{remark}

\section{From width to capacity, and the constant}\label{sec:capacity}

\begin{corollary}[capacity]\label{cor:capacity}
For fixed $\rho>0$ and $Q\ge N$,
\[
\log\Pack_{\rho\sqrt N}\{T_{x,y}\}\;\le\;\frac{2C_S^2\,w_2^2}{\rho^2N}\;\le\;\frac{\kappa}{\rho^2}\,N\log(2Q),
\qquad \kappa:=8C_S^2,
\]
where $C_S$ is the constant in Sudakov's minoration for the process
$X_T=\operatorname{Re}\Tr(G^{*}T)$, whose intrinsic distance is $\Fro{S-T}/\sqrt2$.
\end{corollary}

\begin{proof}
Sudakov's minoration \cite[Thm~7.4.1]{Ver18}, \cite{Tal14}: a family that is $\epsilon$-separated in the
process metric and lies in the index set of a Gaussian process with $\E\sup X_T=w$ has
$\epsilon\sqrt{\log\Pack}\le C_Sw$. At Frobenius separation $\rho\sqrt N$ the process
separation is $\rho\sqrt{N/2}$, and
\[
\E\sup_T\operatorname{Re}\Tr(G^{*}T)\;\le\;w_2,
\]
so inserting Theorem~\ref{thm:width} gives the bound. The hypothesis $Q\ge N$ always holds,
because the target register is part of the register the oracle acts on.
\end{proof}

\begin{callout}{Erratum, stated where a reader will quote it.}
Corollary~\ref{cor:capacity} is \textbf{false in raw Frobenius}, by a factor $N$. It is true
only in the normalized metric of Definition~\ref{def:capacity}. This is a corollary, not a
link in the proof of Theorem~\ref{thm:width}, but it is the statement downstream consumers
quote. Re-check the separation metric before instantiating it.
\end{callout}

\begin{corollary}[qubit lower bound]\label{cor:qubits}
For every $0<\rho<\sqrt2$, $\log\Pack_{\rho\sqrt N}U(N)\ge c\,(1-\rho^2/2)^2N^2$ for an
absolute $c>0$ and all $N\ge3$. Hence no fixed clean two-query architecture realizes every unitary, exactly
or with operator-norm error below $\rho/4$, unless $\log Q=\Omega_\rho(N)$. In qubit count:
\[
\boxed{\;w\;=\;\Omega(2^n)\;}
\]
for a universal clean two-query architecture on $n$ target qubits.
\end{corollary}

\begin{proof}
For fixed $U$ and Haar-random $V$, $\Fro{U-V}^2=2N-2\operatorname{Re}\Tr(U^{*}V)$, so $V$ lies
in the Frobenius ball of radius $\rho\sqrt N$ about $U$ exactly when
$\operatorname{Re}\Tr(U^{*}V)\ge N(1-\rho^2/2)$. The function $V\mapsto\operatorname{Re}\Tr(U^{*}V)$
has Haar mean zero and is $\sqrt N$-Lipschitz in the Frobenius metric, and concentration of
Lipschitz functions on $U(N)$ \cite[Thm~5.17]{Mec19} gives $\Pr(f\ge\E f+t)\le\exp(-c'Nt^2/L^2)$;
at $t=N(1-\rho^2/2)$, $L=\sqrt N$, the ball has Haar measure at most
$\exp(-c'(1-\rho^2/2)^2N^2)$. A maximal $\rho\sqrt N$-separated set is a
$\rho\sqrt N$-covering, so its size is at least the reciprocal of that measure. If an
architecture realizes each member $U_i$ of such a set to operator-norm error below $\rho/4$,
the realized operators are within Frobenius distance $\rho\sqrt N/4$ of the $U_i$ and hence
$\rho\sqrt N/2$-separated; Corollary~\ref{cor:capacity} at scale $\rho/2$ gives
$c(1-\rho^2/2)^2N^2\le4\kappa\rho^{-2}N\log(2Q)$, i.e.\ $\log Q=\Omega_\rho(N)$.
\end{proof}

For comparison, the trivial clean two-query lookup, which stores every entry of $U$ in the
table, uses $\tilde\Theta(N^2)$ qubits, so at this point the threshold sits in
$[\Omega(N),\tilde O(N^2)]$; Theorem~\ref{thm:transport} closes the gap.

\begin{proposition}[the constant matters]\label{prop:constant}
The permutation family of Section~\ref{sec:tight} forces $\kappa\ge\tfrac14$ asymptotically, the
maximum being approached at separation $\rho^2=1$.
\end{proposition}

Computed from Gilbert--Varshamov packings of permutation codes, with $d=N$, in the script
\nolinkurl{verify_crho_requirement.py}:
\begin{center}
\begin{tabular}{lccccc}
\toprule
$d$ & 16 & 256 & 4096 & 65536 & $\to\infty$\\
\midrule
required $\kappa$ & 0.126 & 0.159 & 0.185 & 0.200 & \textbf{0.25}\\
\bottomrule
\end{tabular}
\end{center}
Sudakov's minoration supplies a universal finite constant $C_S$; we do not claim
a numerical upper bound for it. The threshold and tightness results use only a universal
constant. The permutation comparison remains a useful normalization check: any valid
choice of $\kappa$ in Corollary~\ref{cor:capacity} must obey the displayed lower requirement.

\section{Tightness, address length, and the matching upper bound}\label{sec:six}

\subsection{Tight at register dimension of order \texorpdfstring{$N^2$}{N squared}}\label{sec:tight}

\begin{proposition}\label{prop:dlm}
The clean two-query permutation-synthesis algorithm of \cite{DLM26} lies inside
Definition~\ref{def:arch} and attains the ceiling of Corollary~\ref{cor:capacity} up to the
constant factor $4\kappa$.
\end{proposition}

\cite{DLM26} synthesize any permutation unitary $P_\pi$ on $d=N$ dimensions in two clean
queries, querying $\pi$ then $\pi^{-1}$. Index addresses by $(b,x,y)$ with $b\in\{0,1\}$ and
$x,y\in[d]$, so $Q=2d^2$:
\[
g_{0,x,y}=(-1)^{y\cdot\pi(x)},\qquad g_{1,x,y}=(-1)^{x\cdot\pi^{-1}(y)},
\]
\[
C:|x\rangle\mapsto d^{-1/2}\sum_y|0,x,y\rangle,\qquad W=X_b\otimes H_X\otimes H_Y,\qquad
B:|b,x,y\rangle\mapsto\delta_{b,1}\,d^{-1/2}|y\rangle.
\]
$W$ is fixed and independent of $\pi$; a single sign word serves both queries because the
state occupies block $0$ at the first query and block $1$ at the second. Tracing through,
$T_g=P_\pi$ exactly. Verified in \texttt{verify\_permutation\_family.py} over all $24$
permutations at $d=4$ and $200$ random permutations at $d=8$ with residuals below $10^{-14}$.
Since $\Fro{P_\pi-P_\sigma}^2=2\Ham(\pi,\sigma)$, permutation codes at Hamming distance
$\alpha d$ give $\rho^2=2\alpha$, and Gilbert--Varshamov gives
$\log\Pack=\Theta(d\log d)=\Theta(N\log2Q)$.

\begin{corollary}[tight in both factors at $Q=\Theta(N^2)$]\label{cor:tight}
Running Sudakov backwards on the permutation family gives
$w_2\ge c\rho\sqrt d\sqrt{d\log d}=\Theta(N\sqrt{\log2Q})$ at $Q=2N^2$. So at that scale
Theorem~\ref{thm:width} is attained up to a constant, and neither the factor $N$ nor the
factor $\sqrt{\log2Q}$ can be improved by more than a constant there.
\end{corollary}

There is no conflict with the universal lower bound: the permutation family has
packing logarithm $\Theta(N\log N)$ at fixed separation, whereas covering all
of $U(N)$ requires $\Theta(N^2)$. It establishes tightness of the capacity bound,
not universality at $Q=\Theta(N^2)$.

Padding an architecture with unused workspace raises $\log Q$ without changing $w_2$, so
this is a statement about the scale $Q=\Theta(N^2)$, not about every polynomial. At
exponential $Q$ the logarithm is necessary as well: by Theorem~\ref{thm:transport}, at
$\log Q=\Theta(N\log(1/\varepsilon))$ the class $\varepsilon$-covers $U(N)$ in operator norm,
and since $|\Tr(G^{*}(T-U))|\le\|G\|_{S_1}\|T-U\|_{\op}$ this gives
$w_2\ge(1-\varepsilon)\,\E\|G\|_{S_1}=\Theta(N^{3/2})$ there. No bound of the form
$N\cdot\polylog(N)$ can therefore hold uniformly in $Q$.

\subsection{The oracle's input length alone}\label{sec:address}

Theorem~\ref{thm:width} bounds the total register dimension. The next theorem removes the
workspace from the accounting entirely. In this section and the next, $M$ denotes a number
of addresses; the coefficient matrix $M$ of Section~\ref{sec:width} does not appear.

\begin{theorem}[address-length lower bound]\label{thm:address}
Let the oracle act on the $Q$-dimensional register through $M$ addresses of arbitrary rank,
\[
D_x=\sum_{j=1}^{M}x_j\,P_j,\qquad x\in\T^M,
\]
with $P_1,\dots,P_M$ mutually orthogonal projections summing to $I_Q$ (for a Boolean table on
$\ell$ input bits, $M=2^\ell$ and $P_j$ projects onto address $j$ tensored with the whole
workspace). Then for every two-query architecture
\[
\E_G\sup_{x,y\in\T^M}\bigl|\Tr(G^{*}BD_xWD_yC)\bigr|\;\le\;2\,\Fro{B}\Fro{C}\sqrt{\log(2MN)}
\;\le\;2N\sqrt{\log(2MN)},
\]
and consequently universal clean two-query synthesis, exact or to operator-norm error below
an absolute constant, requires $\log M=\Omega(N)$, i.e.\ an oracle input length of
$\Omega(2^n)$ bits, regardless of the number of ancillas.
\end{theorem}

\begin{proof}
Let $L:=\sum_j\ran(P_jB^{*})$ and $R:=\sum_j\ran(P_jC)$, subspaces of $\C^Q$ of dimension at
most $MN$ each. Because the $P_j$ are mutually orthogonal, $L$ is the orthogonal direct sum
of the $\ran(P_jB^{*})$ and is invariant under every $P_j$; likewise $R$. Let $V_L,V_R$ be
isometries onto $L,R$ and $\Pi_L=V_LV_L^{*}$, $\Pi_R=V_RV_R^{*}$. Since $(BP_j)^{*}=P_jB^{*}$
has range in $L$, $BP_j=BP_j\Pi_L$, and summing over $j$ gives $B=B\Pi_L$; similarly
$C=\Pi_RC$. Hence
\[
BD_xWD_yC\;=\;(BV_L)\Bigl(\sum_jx_j\,V_L^{*}P_jV_L\Bigr)\bigl(V_L^{*}WV_R\bigr)\Bigl(\sum_ky_k\,V_R^{*}P_kV_R\Bigr)(V_R^{*}C).
\]
In an orthonormal basis of $L$ adapted to its decomposition, $\sum_jx_jV_L^{*}P_jV_L$ is
diagonal with each $x_j$ repeated $\rank(P_j|_L)$ times; it is a unimodular diagonal, and
likewise on $R$. The compressed middle $V_L^{*}WV_R$ is a contraction, and
$\Fro{BV_L}^2=\Tr(B^{*}B\Pi_L)=\Fro{B}^2$, $\Fro{V_R^{*}C}=\Fro{C}$. Theorem~\ref{thm:width}
applies to the compressed architecture: its statement has a square middle layer, but its
proof never uses squareness (Remark~\ref{rem:zero}), and with a $\dim L\times\dim R$ middle
the Khintchine step gives $\sqrt{\log(\dim L+\dim R)}\le\sqrt{\log(2MN)}$. The supremum over
all unimodular diagonals on $L$ and $R$ dominates the supremum over the repeated patterns.
Corollary~\ref{cor:capacity} then reads $\log\Pack\le\kappa\rho^{-2}N\log(2MN)$, and the
$\Omega(N^2)$ demand of a Haar-scale code, with the error tolerance of
Corollary~\ref{cor:qubits}, forces $\log M=\Omega(N)$.
\end{proof}

For a bit-flip oracle $O_f|x,b\rangle=|x,b\oplus f(x)\rangle$ the eigen-decomposition of the
response bit gives $M+1$ addresses; nothing changes. The two selectors may also use
different address decompositions, as for two tables of lengths $M_1$ and $M_2$: the proof is
unchanged with $\dim L\le M_1N$ and $\dim R\le M_2N$. The factorization, the dimension count,
the preserved norms and the contractivity are checked on eighty random architectures with
address projectors of arbitrary rank, and with two different decompositions, in
\texttt{verify\_address\_compression.py}.

\subsection{The matching upper bound}\label{sec:upper}

\begin{theorem}[permutation transport]\label{thm:transport}
For every $\varepsilon\in(0,\tfrac1{10})$ there are a target-independent encoder and decoder
such that every $U\in U(N)$ is implemented by a clean two-query circuit with
\begin{center}
\begin{tabular}{ll}
\toprule
resource & cost\\
\midrule
Boolean tables & one\\
Boolean queries & two\\
oracle input length & $O(N\log(1/\varepsilon))$ bits\\
workspace & $O(N\log(1/\varepsilon))$ qubits\\
ordinary gates & $\poly(N,1/\varepsilon)$\\
half-diamond error, all reference-entangled inputs & at most $\varepsilon$\\
extra target-dependent advice, postselection & none\\
\bottomrule
\end{tabular}
\end{center}
The construction restores its workspace to within the stated error, so it is clean; the
clean-output compression $T_g$ is within operator-norm distance $3\varepsilon/4$ of
$U$; and the circuit is of the repeated-sign-word form $T_g$ of Definition~\ref{def:arch},
one $\pm1$ table serving both queries.
\end{theorem}

\paragraph{Construction, condensed.} The full statement with all estimates is in the
Appendix~\ref{app:transport}, also available as
\nolinkurl{notes/permutation-transport.md} in the repository.
\nolinkurl{verify_koopman_transport.py} checks the ingredients at small $N$.

Encode a state as a function on a finite grid rather than as a list of columns. Quantize the
standard normal distribution into $J$ equiprobable intervals with conditional means $q_a$,
put $s^2=\tfrac1J\sum_aq_a^2$ and $e^2=1-s^2$; a direct estimate gives $e^2\le32J^{-1/3}$
(natural logarithms throughout), so $J=\poly(1/\varepsilon)$ makes $e$ as small as needed,
independently of $N$. For addresses $a\in[J]^{2N}$, $M=J^{2N}$, define $w_a\in\C^N$ by
$(w_a)_k=(q_{a_k}+iq_{a_{N+k}})/(s\sqrt2)$ and the fixed encoding $(E\psi)_a=w_a^{*}\psi/\sqrt M$.
Independence, zero means and the normalization give $E^{*}E=I_N$ exactly.

For a target $U$ with realification $O_U$, draw $X\sim\mathcal N(0,I_{2N})$ and quantize
coordinatewise with the cell map $\theta$: $\alpha=\theta(X)$, $\beta=\theta(O_UX)$. Both are
uniform on $[J]^{2N}$, so $\Xi_{ba}=M\Pr(\alpha=a,\beta=b)$ is doubly stochastic and, by
Birkhoff--von Neumann \cite{Bir46}, a convex combination $\sum_{\ell=1}^{m}p_\ell P_{\pi_\ell}$
of at most $m\le M^2$ permutation matrices. The quantization residual $r(X)=X-q_{\theta(X)}$
has covariance $e^2I_{2N}$, and $q_{\theta(O_UX)}-O_Uq_{\theta(X)}=O_Ur(X)-r(O_UX)$, whence,
using
$(a-b)(a-b)^{\top}\preceq2aa^{\top}+2bb^{\top}$,
\[
\sum_\ell p_\ell\,\Delta_{\pi_\ell}^{*}\Delta_{\pi_\ell}\;\preceq\;\frac{4e^2}{s^2}\,I_N,
\qquad\Delta_\pi:=P_\pi E-EU.
\]
This is a uniform operator bound, valid for every input. To make the mixture a single
permutation, put $\delta:=\varepsilon/2$, take $J$ with $e^2\le\delta^2/100$, and add a
uniform coherent label register over $R\ge100M^2/\delta^2$ copies, the least power of two satisfying this bound, fixed by
$M$ and $\delta$ alone and hence independent of the target, with integer multiplicities
$k_\ell\approx Rp_\ell$, $\sum_\ell|p_\ell-k_\ell/R|\le2m/R$; set
$\Pi_U(j,a)=(j,\pi_{\ell(j)}(a))$ with the enlarged, still target-independent, encoding
$\tilde E\psi=|+_R\rangle\otimes E\psi$. Then
\[
\bigl\|\Pi_U\tilde E-\tilde EU\bigr\|_{\op}^2
=\Bigl\|\sum_\ell\tfrac{k_\ell}{R}\Delta_{\pi_\ell}^{*}\Delta_{\pi_\ell}\Bigr\|_{\op}
\;\le\;\frac{4e^2}{s^2}+\frac{8m}{R}\;<\;\delta^2,
\]
so the label register returns to $|+_R\rangle$ up to $\delta$, uncorrelated with the input.
The permutation $\Pi_U$ and its inverse are stored in one Boolean table on $1+2t$ input
bits, $t=\log_2(MR)=O(N\log(1/\varepsilon))$, and are applied by two queries in the standard
way (load $\Pi_U(z)$ into a blank register, swap, unload with the inverse table). The fixed
encoder is an explicit polynomial-size circuit and its inverse decodes. Total isometry error
at most $3\varepsilon/4$: $\delta=\varepsilon/2$ from the label register and $\varepsilon/8$
each for the encoder and the decoder. Two isometries at operator distance $\gamma$ induce
channels at half-diamond distance at most $\gamma$. The full estimates are in the note.

The lower bound of Theorem~\ref{thm:address} shows that for this mechanism, and indeed for
any monomial action between fixed encodings, the exponential address length is unavoidable;
the note records that converse and its one loophole, a target-dependent final ancilla
state, which does not affect the clean model treated here.

\begin{corollary}[the thresholds]\label{cor:thresholds}
For clean two-query synthesis of all $n$-qubit unitaries at operator-norm error
$\varepsilon$ below an absolute constant, both the qubit count $w^{\star}$ and the oracle
input length $\ell^{\star}$ satisfy
\[
\Omega(2^n)\;\le\;w^{\star},\;\ell^{\star}\;\le\;O\bigl(2^n\log(1/\varepsilon)\bigr).
\]
\end{corollary}

\section{The garbage model: what the two-query problem costs there}\label{sec:garbage}

In the garbage model the workspace may end in a target-dependent state $|j_g\rangle$,
independent of the input, or may simply be discarded. \cite{LMW24} and \cite{DLM26} state
their one-query bounds there. This section gives a sufficient route from an
insertion-class capacity bound to a lower bound in that model.

\paragraph{Setting.} A $t$-query circuit
$V(g)=A_t(D_g\otimes I_k)A_{t-1}\cdots(D_g\otimes I_k)A_0$ on $\C^Q\otimes\C^k$, where $Q$ is the
number of addresses and $k$ the workspace dimension, restricted to the clean input,
$S_g:=V(g)\iota$ with $\iota\psi=\psi\otimes|0\rangle$, is an isometry
$\C^N\to\C^Q\otimes\C^k\cong\C^N\otimes\C^D$, with $D=Qk/N$ the junk dimension, whose entries are
multilinear polynomials of degree at most $t$ in the signs. In the \textbf{keep form}
$S_g\psi=(U_g\psi)\otimes|j_g\rangle$. The \textbf{discard form} only requires the induced
channel to be within diamond distance $\varepsilon$ of $U_g$; since a unitary channel has
Kraus rank one, continuity of the Stinespring representation recovers the keep form up to
$O(\sqrt\varepsilon)$ in operator norm: \cite{KSW08} gives the required unitary on a common
enlarged environment, which suffices here because embedding the junk register isometrically
in a larger one changes none of the conjugated objects below, and \cite{vE23} gives it on
the junk register itself with constant $\sqrt2$. So everything below transfers to the
discard form at constant error.

\begin{lemma}[junk-blindness]\label{lem:junk}
Call a functional $L$ on isometries junk-blind if $L(S)=L((I\otimes W)S)$ for every unitary
$W$ on the junk register, equivalently if it depends only on the induced channel. If $L$ is
junk-blind and linear then $L\equiv0$.
\end{lemma}

\begin{proof}
Take $W=e^{i\theta}I$: junk-blindness gives $L(S)=L(e^{i\theta}S)=e^{i\theta}L(S)$ for all
$\theta$; for a real-linear $L$ take $\theta=\pi$.
\end{proof}

The lemma rules out a nonzero junk-blind linear functional of the amplitude operator itself.
Conjugation provides a useful quadratic invariant with an explicit $2t$-insertion
representation. This proves sufficiency of the transfer below, not necessity of
oracle degree $2t$: substitution of an oracle circuit into an invariant can cause
degree cancellations. For example, $D_gD_g=I$ is a two-query product with constant
conjugated observables. No lower bound on the degree of every possible invariant,
or on every width method, is asserted.

\begin{theorem}[conjugation transfer]\label{thm:transfer}
For every fixed compatible tuple of contractions $(Y_0,\ldots,Y_{2t})$, consider the
$N\times N$ family $\{Y_{2t}D_gY_{2t-1}\cdots D_gY_0:g\in\{\pm1\}^Q\}$ with
$D_g=\sum_{j\le Q}g_jP_j$ for mutually orthogonal projectors $P_j$ of arbitrary rank summing
to the identity of a register of arbitrary dimension, as in Theorem~\ref{thm:address}.
Suppose each such family satisfies a capacity bound
$\log\Pack_{\rho\sqrt N}\le C_\rho N\,\phi_N(Q)$, uniformly over the fixed contractions,
projectors and workspace dimension. Packing is over $g$ alone, not over these fixed
architectural choices. Here $\phi_N$ is a rate function of the address count, allowed to
depend on $N$ but not on workspace dimension.
Then no universal garbage-$t$ synthesizer exists with
$\phi_N(Q)=o(N)$, for every workspace and junk dimension, with no change to the oracle, at
sufficiently small constant diamond error. At $t=1$, Theorem~\ref{thm:address} gives $\phi_N(Q)=\log(2QN)$.
\end{theorem}

\begin{proof}
Fix a balanced reflection $R$ ($R=R^{*}$, $R^2=I$, $\Tr R=0$) and set
$T_g:=S_g^{*}(R\otimes I)S_g$. Because $D_g\otimes I$ is Hermitian, expanding $S_g^{*}$ by
reversal writes $T_g$ exactly as a $2t$-insertion product of the same sign word on
$\C^Q\otimes\C^k$, with address projectors $P_j\otimes I_k$ and contractions in every slot (the
middle slot is $A_t^{*}(R\otimes I)A_t$; the product is palindromic,
$T_g=Y_g^{*}\,A_t^{*}(R\otimes I)A_t\,Y_g$ with $Y_g=(D_g\otimes I)A_{t-1}\cdots(D_g\otimes I)A_0\iota$).
In the keep form $T_g=\langle j_g|j_g\rangle\,U_g^{*}RU_g=U_g^{*}RU_g$: the junk cancels by
the normalization of a unit vector, with no reference vector, no alignment and no
dependence on the junk dimension. The map $U\mapsto U^{*}RU$ lands in the balanced
reflections, a copy of the Grassmannian $\Gr(N/2,N)$ of real dimension $N^2/2$ on which every
point has Frobenius norm $\sqrt N$. Here is a direct packing justification:
for Haar $V$, $f(V)=\Tr(RV^*RV)$ is real, has mean zero, and is
$2\sqrt N$-Lipschitz in Frobenius distance. A ball of radius
$\rho\sqrt N$ about $R$ is the event $f(V)\ge N(1-\rho^2/2)$.
The same concentration inequality as in Corollary~\ref{cor:qubits} bounds
its measure by $\exp(-c_\rho N^2)$ for fixed $\rho<\sqrt2$.
Invariance gives the same measure at every orbit point; a maximal separated
set covers the orbit and therefore has at least $e^{c_\rho N^2}$ points,
each of the form $U_i^*RU_i$. A universal synthesizer must synthesize $U_1,\dots,U_m$, so the
single family $\{T_g\}$ contains an $e^{c_\rho N^2}$-point separated code, contradicting the
assumed capacity bound when $\phi_N(Q)=o(N)$. For the discard form no
Stinespring continuity is needed: if
$\|\Phi_g-\Ad_U\|_\diamond\le\varepsilon$, channel duality gives
$\|\Phi_g^*(R)-U^*RU\|_{\op}\le\varepsilon$.
Thus the normalized Frobenius error is at most $\varepsilon$,
absorbed by taking a sufficiently small constant service error.
\end{proof}

The reversal identity, the junk cancellation and the orbit properties are checked in
\texttt{verify\_conjugation\_transfer.py}.

\begin{corollary}[one query, garbage model]\label{cor:onequery}
At $t=1$ the hypothesis of Theorem~\ref{thm:transfer} is Theorem~\ref{thm:address}, applied
to two insertions of a repeated real sign word. Hence no one-query circuit, with garbage
allowed and with any number of ancillas, synthesizes every unitary at constant diamond error
unless its oracle input length is $\Omega(2^n)$.
\end{corollary}

This recovers, by a different route and at constant error only, the address-length form of
the one-query bound of \cite{LMW24}; we claim no improvement on their statement, which also
rules out vanishing advantage.

\begin{remark}[what is not claimed]\label{rem:notclaimed}
At $t=2$ the hypothesis of Theorem~\ref{thm:transfer} is a clean \textbf{four}-insertion
width bound, which is open. By the palindromic form it suffices to bound the sandwich
subclass $Y_g^{*}XY_g$ with $Y_g$ a two-insertion product whose left endpoint is the identity
of the full register: exactly the case in which the endpoint budget that makes step~3 of
Theorem~\ref{thm:width} work is unavailable, since that identity has squared Frobenius norm
equal to the register dimension. So the unrestricted garbage-model two-query problem is
open, and this paper does not claim it. A restricted two-query garbage bound, valid when the
junk register has fewer than $2n-O(1)$ qubits, exists in the author's working records but
has not passed an independent audit and is not included.
\end{remark}

\section{Scope, and depth three and beyond}\label{sec:scope}

\paragraph{What Theorem~\ref{thm:width} covers.} Exact and approximate clean two-query
synthesis; arbitrary contractive $B,W,C$; real sign oracles, diagonal-unitary oracles and two
different tables; ancillas explicitly inside $\log(2Q)$, and, via Theorem~\ref{thm:address},
the oracle input length alone.

\paragraph{What it does not cover.}
\begin{enumerate}[leftmargin=*]
\item \textbf{Three or more queries.} Nothing here bears on $t\ge3$. The natural conjecture,
that the width of the $k$-query clean class is $O(N\cdot k\cdot\polylog Q)$, is pinned in
exponent by \cite{Ros26}: every unitary is implementable at $k=\tilde O(\sqrt N)$ with
$\log Q=\poly(n)$, so the width must reach the $\Theta(N^{3/2})$ width of the unitary group
there, which forces at least linear growth in $k$. A proof of the conjecture would resolve
the Aaronson--Kuperberg question negatively in the clean model. Less would do: any bound of the form
$N\cdot\poly(k,\log Q)$ already excludes universal clean synthesis at polynomial depth and
polynomial qubit count, since such a class could not $\varepsilon$-cover $U(N)$, whose width
is $\Theta(N^{3/2})$ (the argument of Corollary~\ref{cor:tight}); and by
Theorem~\ref{thm:transfer} the same bound at $2k$ insertions, in address form, would exclude
it in the garbage model as well. The two-query proof does not extend as it stands: with a
third insertion the coefficient matrix depends on the selector.
\item \textbf{Parallel versus adaptive.} \cite{LMW24} rule out polynomially many
\emph{parallel} queries. The difficulty at $t\ge3$ is therefore entirely in interleaving.
\item \textbf{One query.} \cite[Thm~6.7]{AK07} is stronger at $t=1$ in the exact model: $4^N$
with at most $2N$ relevant table entries. This counts Boolean variables, not oracle
input bits. Banerjee \cite{Ban25} sharpens the matrix count to $2^N$ under the
exact every-oracle-clean assumption.
\item \textbf{Completely bounded forms.} \cite[Thm~1.3]{ABP18} characterizes scalar expected
outputs of $t$-query algorithms by completely bounded forms of degree $2t$.
Our matrix-valued amplitude operator is separately multilinear of degree $t$.
These are related factorizations, but the cited scalar theorem is not an equivalence
with the full matrix-valued unit ball. Such an identification is not used or claimed here.
\end{enumerate}

\paragraph{Consistency checks.} The bound does not forbid the trivial clean two-query lookup
($w=\Theta(N^2)$), is attained up to $\log(1/\varepsilon)$ by Theorem~\ref{thm:transport}, and
does not contradict \cite{Ros26}, whose algorithm uses $t=O(\sqrt N)$ queries. Any extension
of this method yielding $t=\omega(\sqrt N)$ in the clean model would contradict \cite{Ros26}
and be in error.

\section{Reproducibility}\label{sec:repro}

\texttt{verify/run\_all.py} in the repository runs nine deterministic scripts
(standard library and numpy),
each exiting nonzero on any failed check: Theorem~\ref{thm:width} in both forms (exact
expansion, exhaustive absorption on finite sign cubes and checks at sampled or
alternating-maximization phase vectors, the closed-form variance
operators against Monte Carlo and their domination), Proposition~\ref{prop:unitarity},
Propositions~\ref{prop:invisible} and \ref{prop:pair}, Proposition~\ref{prop:constant},
Proposition~\ref{prop:dlm}, the factorization of Theorem~\ref{thm:address}, the ingredients of
Theorem~\ref{thm:transport} at small $N$, and the identities of Section~\ref{sec:garbage}.
The \texttt{lean/} directory holds a kernel-checked formalization of the deterministic inequalities
of Corollary~\ref{cor:real}, with the analytic inputs as explicit hypotheses; its
\texttt{README} states exactly what is and is not certified.

\section*{Acknowledgements and disclosure}

\paragraph{Use of language models.}
Large language models were used extensively throughout this work for proof discovery,
implementation, verification engineering, adversarial auditing, and editorial assistance.
The workflow used separate, independently instantiated sessions, including fresh-context
auditing sessions instructed to refute rather than confirm. The author directed the
research program, set the verification standards, made the promotion, correction, and
publication decisions, and assumes sole responsibility for this work. These AI audits
are not external peer review.

The proofs are analytic, model-assisted arguments; exact computation and repeated audits
do not by themselves validate their analytic dependencies. The scripts check specified
identities and finite instances, using exact rational arithmetic for designated controls
and floating-point calculations elsewhere; numerical searches do not certify global
optima. The Lean kernel checks seven selected declarations, including conditional
assembly results whose analytic inputs remain explicit hypotheses. The complete width,
packing, address-compression, and transport proofs are not formalized in Lean.

\appendix
\section{Complete permutation-transport construction}\label{app:transport}

This appendix proves Theorem~\ref{thm:transport} in the fixed clean-isometry
metric, including the finite encoding, rounding, circuit implementation and
all reference-entangled inputs. No efficient algorithm for constructing the
target-dependent Boolean truth table is assumed or claimed.

\subsection{Quantization with dimension-independent error}
Let $J\ge4$ be a power of two. Partition a real standard normal variable $X$
into $J$ equiprobable intervals, with cell index $\theta(X)$ and conditional
means $q_a=\E[X\mid\theta(X)=a]$. Set
\[
 s^2=J^{-1}\sum_a q_a^2,\qquad e^2=\E(X-q_{\theta(X)})^2=1-s^2.
\]
The means sum to zero. For completeness, $e^2\le32J^{-1/3}$:
write $g(u)=\Phi^{-1}(u)$ and clip it to $[-T,T]$ to obtain $g_T$,
which is Lipschitz with constant $L_T=\sqrt{2\pi}e^{T^2/2}$.
Orthogonal projection onto cell-constant functions is the best $L^2$
approximation, so $e\le\|g-g_T\|_2+L_T/J$. For $T\ge1$,
\[
 \|g-g_T\|_2^2\le \E[X^2{\bf1}_{|X|>T}]
 =2(T\phi(T)+1-\Phi(T))\le2(T+1)\phi(T).
\]
With $T=\sqrt{\log J}$ this gives
\[
 e^2\le\frac{4(\sqrt{\log J}+1)}{\sqrt{2\pi J}}+\frac{4\pi}{J}
 \le32J^{-1/3}.
\]
One elementary check of the last bound is to multiply by $J^{1/3}$:
$(\sqrt u+1)e^{-u/6}\le\sqrt3 e^{-1/2}+1<3$ for $u\ge0$
(maximize $\sqrt u e^{-u/6}$ at $u=3$), and
$4\pi J^{-2/3}\le4\pi/4^{2/3}<5$.
Thus the displayed constant is conservative.

\subsection{A fixed finite isometry and its circuit}
Put $\Omega=[J]^{2N}$, $M=J^{2N}$, and
\[
 (w_a)_k=\frac{q_{a_k}+i q_{a_{N+k}}}{s\sqrt2},
 \qquad (E\psi)_a=\frac{w_a^*\psi}{\sqrt M}.
\]
Independence and zero means give $M^{-1}\sum_a w_aw_a^*=I_N$,
so $E^*E=I_N$ exactly. Define orthonormal coordinate states
\[
 |g\rangle=J^{-1/2}\sum_a|a\rangle,\qquad
 |h\rangle=(s\sqrt J)^{-1}\sum_a q_a|a\rangle.
\]
Then $E|j\rangle$ is the superposition of the configuration with $h$
in coordinate $j$ and the one with $h$ in coordinate $N+j$, with
coefficients $1/\sqrt2$ and $-i/\sqrt2$; all other coordinates hold $g$.
Coherently convert the binary label $j$ into these two possible single
excitation positions, and erase the binary label by computing it back
from the occupied position. Apply on each coordinate a fixed unitary
taking $|0\rangle$ to $g$ and $|1\rangle$ to $h$.
This gives an explicit unitary extension of $E$ with
$O(N(n+1)^2+NJ^2\poly(\log J,\log(N/\varepsilon)))$ gates after
finite-precision compilation. Compile the whole encoder to operator
error $\varepsilon/8$ and use its circuit inverse for decoding.
Every angle is target independent. This gate count is polynomial in
$N,1/\varepsilon$, not in $n$.

\subsection{Transport coupling and its operator error}
For a fixed target $U$, let
\[
 O_U=\begin{pmatrix}\Re U&-\Im U\\ \Im U&\Re U\end{pmatrix}.
\]
Take $X$ standard normal in $\R^{2N}$ and set
$\alpha=\theta(X)$, $\beta=\theta(O_UX)$ coordinatewise.
Both marginals are uniform on $\Omega$, so
$\Xi_{ba}=M\Pr(\alpha=a,\beta=b)$ is doubly stochastic.
It is a convex combination $\sum_{\ell=1}^{m}p_\ell P_{\pi_\ell}$
with $m\le M^2$. To see finiteness directly, Hall's condition follows
from equal row and column sums on the positive support; subtract the
smallest matching entry at each step. Row and column sums stay equal,
at least one positive entry disappears, and the process ends in at most
$M^2$ steps. The weights sum to one.

Write $r(X)=X-q_{\theta(X)}$. Its coordinates are independent,
centered, and have variance $e^2$, and the same distributional statement
holds for $r(O_UX)$. Therefore
\[
 \E[(q_\beta-O_Uq_\alpha)(q_\beta-O_Uq_\alpha)^\top]
 \preceq4e^2 I_{2N}.
\]
Here we used $q_\beta-O_Uq_\alpha=O_Ur(X)-r(O_UX)$ and
$(a-b)(a-b)^\top\preceq2aa^\top+2bb^\top$; no independence
between these two residual vectors is needed. For real $x,y$, the map
$(x,y)\mapsto(x+iy)/(s\sqrt2)$ has squared real-linear operator norm $1/(2s^2)$.
To transport the covariance as a complex Hermitian matrix, however, use its
complex-linear extension $L=[I_N\;\;iI_N]/(s\sqrt2)$, for which
$LL^*=I_N/s^2$. Thus, writing $\Sigma\preceq4e^2I_{2N}$ for the real
covariance above, $L\Sigma L^*\preceq4e^2LL^*$ gives
\[
 C_U:=\E[(w_\beta-Uw_\alpha)(w_\beta-Uw_\alpha)^*]
 \preceq(4e^2/s^2)I_N.
\]
For $\Delta_\pi=P_\pi E-EU$, expansion at output coordinate
$b=\pi(a)$ gives, for every $\psi$,
\[
 \sum_\ell p_\ell\|\Delta_{\pi_\ell}\psi\|^2
 =\E|w_\alpha^*\psi-w_\beta^*U\psi|^2
 =\psi^*U^*C_UU\psi.
\]
Consequently
\begin{equation}\label{eq:transport-covariance}
 \sum_\ell p_\ell\Delta_{\pi_\ell}^*\Delta_{\pi_\ell}
 \preceq (4e^2/s^2)I_N .
\end{equation}
This bounds all inputs simultaneously, not only the maximally mixed input.

\subsection{One coherent permutation and a fixed label state}
Set $\delta=\varepsilon/2$ and choose $J$ with $e^2\le\delta^2/100$.
Choose the least power of two $R\ge100M^2/\delta^2$ depending only on
$M,\delta$. Round the weights to integers $k_\ell\ge0$ with
$\sum k_\ell=R$ and $\sum|p_\ell-k_\ell/R|\le2m/R$:
round down first, then distribute the remaining slots.
List $k_\ell$ copies of each permutation and define
\[
 \Pi_U(j,a)=(j,\pi_{\ell(j)}(a)),\qquad
 \widetilde E\psi=|+_R\rangle\otimes E\psi.
\]
Orthogonality of the label states gives the exact identity
\[
 \|\Pi_U\widetilde E-\widetilde EU\|_{\op}^2
 =\left\|\sum_\ell(k_\ell/R)\Delta_{\pi_\ell}^*\Delta_{\pi_\ell}\right\|_{\op}.
\]
Since $\|\Delta_\pi\|\le2$, \eqref{eq:transport-covariance} implies
\[
 \|\Pi_U\widetilde E-\widetilde EU\|_{\op}^2
 \le4e^2/s^2+8m/R<\delta^2.
\]
Thus the label returns approximately to the same uniform state;
it is not measured, discarded, or allowed to retain target-dependent junk.
The binary coordinate and label registers have
$t=\log_2(MR)=O(N\log(1/\varepsilon))$ qubits in total.

\subsection{Two Boolean queries, complete address count, and error}
On a direction bit $d$ and two $t$-bit strings $z,y$, define one table
\[
 f_U(d,z,y)=
 \begin{cases}
  \langle\Pi_U(z),y\rangle_{\mathbb F_2},&d=0,\\
  \langle\Pi_U^{-1}(z),y\rangle_{\mathbb F_2},&d=1.
 \end{cases}
\]
The full oracle input length is $1+2t$. A phase query sandwiched
between Hadamards on $y$ implements XOR lookup of the entire
$t$-bit answer, by character orthogonality. Forward lookup, register
swap, and inverse lookup take $|z,0\rangle$ to $|\Pi_U(z),0\rangle$.
The direction bit and the phase-answer qubit are restored by fixed
gates. Both lookups query the same table, in two disjoint tagged sections.
The remaining gates and encoder/decoder are fixed before $U$ is chosen.

Decoding the transport error and including both compilation errors gives
full-isometry distance at most $\delta+2(\varepsilon/8)=3\varepsilon/4$
from $\psi\mapsto U\psi\otimes|0\cdots0\rangle$.
Tensoring with a reference preserves this operator bound.
The half trace distance between the corresponding pure states is at most
their vector distance; partial trace and convexity extend the guarantee
to arbitrary mixed reference-entangled inputs. Thus the half-diamond
error is at most $\varepsilon$, with approximate cleanup in the stronger
full-isometry sense. No bound on the offline time or size of the truth
table is claimed.

\subsection{A converse for monomial action between fixed encodings}
Let $E,F:\C^N\to\C^D$ be fixed isometries and let the only
target-dependent operation be a monomial $P$.
Write $e_i=E^*|i\rangle$, $f_j=F^*|j\rangle$, and define
\[
 a(U)=\max_{i,j:e_i,f_j\ne0}
 \frac{|\langle Ue_i,f_j\rangle|}{\|e_i\|\|f_j\|}.
\]
For every permutation and every choice of its phases,
$|\Tr(U^*F^*PE)|\le N a(U)$ by Cauchy--Schwarz and
$\sum\|e_i\|^2=\sum\|f_j\|^2=N$.
If the two isometric channels are within half-diamond error $\varepsilon$,
their pure normalized Choi states have squared overlap at least
$1-\varepsilon^2$, forcing $a(U)^2\ge1-\varepsilon^2$.
For fixed unit vectors $v,w$ and Haar $U$,
$\Pr(|\langle Uv,w\rangle|^2\ge1-\varepsilon^2)=\varepsilon^{2(N-1)}$.
A union bound over at most $D^2$ pairs gives universality only if
$D^2\varepsilon^{2(N-1)}\ge1$, or
$\log_2D\ge(N-1)\log_2(1/\varepsilon)$.
This is a converse for fixed output encoding including its ancilla,
not for a target-dependent final garbage state.

\section{Exact finite control for Proposition~\ref{prop:pair}}\label{app:pair}
This example counts scalar unitary matrices, not distinct channels:
at $N=1$ all these matrices induce the same channel. It is not used
in any lower bound.
Index the eight coordinates by $x\in\mathbb F_2^3$ in lexicographic
order. Write $\chi_a(x)=(-1)^{a\cdot x}$,
$A(a_0,a_1,a_2)=(a_2,a_0+a_2,a_1)$ and $f=8^{-1/2}(1,\ldots,1)$.
Both $A$ and $I+A$ are invertible over $\mathbb F_2$.
Take $C=f$, $B=f^*$ and define
\[
 W_1=\frac18\sum_a\chi_{Aa}\chi_a^\top,\qquad
 W_2(\lambda)=\frac18\sum_a\lambda_a\chi_a\chi_{(I+A)a}^\top.
\]
The character bases are orthogonal, so both mixers are unitary when
$|\lambda_a|=1$. For $g=\pm\chi_a$, direct substitution gives
$T_g=\pm\lambda_a$ and no leakage.
Use $\lambda_a=1$ for the first architecture. For the second, enumerate
$a$ by $j=0,\ldots,7$ and set
$\lambda_a=(1-j^2+2ij)/(1+j^2)$.
The sixteen values $\pm\lambda_a$ are distinct.
For completeness, the exact finite certificate in
\texttt{verify/verify\_exact\_controls.py} enumerates all $256$ sign words,
forms $v=D_gW_1g$, and tests whether all entries of $D_gW_2v$ are
identical. In rational real and imaginary arithmetic the accepted words
are precisely the sixteen $\pm\chi_a$ for both architectures.
This is an exhaustive finite certificate, with no floating-point
tolerance. The frames before $W_2$ agree because $C,W_1$ agree.

\begin{thebibliography}{BBCMW98}

\bibitem[AK07]{AK07} S.~Aaronson, G.~Kuperberg. Quantum versus classical proofs and advice.
\emph{Theory of Computing} 3(7):129--157, 2007. arXiv:quant-ph/0604056.

\bibitem[LMW24]{LMW24} A.~Lombardi, F.~Ma, J.~Wright. A one-query lower bound for unitary
synthesis and breaking quantum cryptography. \emph{STOC 2024}, 979--990. arXiv:2310.08870.

\bibitem[DLM26]{DLM26} F.~Dong, A.~Lombardi, F.~Ma. Explicit separations for one-query unitary
synthesis. arXiv:2607.26478, 2026.

\bibitem[BY26]{BY26} Z.~Brakerski, H.~Yuen. On scalable pseudorandom unitaries and the unitary
synthesis problem. arXiv:2605.09957, 2026.

\bibitem[Ros26]{Ros26} G.~Rosenthal. Query and depth upper bounds for quantum unitaries via
Grover search. \emph{Quantum}, 2026, article 2144. arXiv:2111.07992.

\bibitem[ABP18]{ABP18} S.~Arunachalam, J.~Bri\"et, C.~Palazuelos. Quantum query algorithms are
completely bounded forms. \emph{ITCS 2018}; \emph{SIAM J. Comput.} 48(3):903--925, 2019.
arXiv:1711.07285.

\bibitem[KSW08]{KSW08} D.~Kretschmann, D.~Schlingemann, R.~F.~Werner. The
information-disturbance tradeoff and the continuity of Stinespring's representation.
\emph{IEEE Trans. Inf. Theory} 54(4):1708--1717, 2008.

\bibitem[vE23]{vE23} F.~vom Ende. Progress on the Kretschmann--Schlingemann--Werner
conjecture. arXiv:2308.15389, 2023.

\bibitem[Tro15]{Tro15} J.~A.~Tropp. An introduction to matrix concentration inequalities.
\emph{Foundations and Trends in Machine Learning} 8(1--2):1--230, 2015.

\bibitem[Tal14]{Tal14} M.~Talagrand. \emph{Upper and lower bounds for stochastic processes.}
Ergebnisse der Mathematik 60, Springer, 2014.

\bibitem[Ver18]{Ver18} R.~Vershynin. \emph{High-Dimensional Probability.} Cambridge University
Press, 2018. (Sudakov--Fernique, Thm~7.2.11; Sudakov minoration, Thm~7.4.1.)

\bibitem[Mec19]{Mec19} E.~S.~Meckes. \emph{The Random Matrix Theory of the Classical Compact
Groups.} Cambridge University Press, 2019. (Concentration of Lipschitz functions on $U(N)$,
Thm~5.17.)

\bibitem[BBCMW98]{BBCMW98} R.~Beals, H.~Buhrman, R.~Cleve, M.~Mosca, R.~de~Wolf. Quantum lower
bounds by polynomials. \emph{FOCS 1998}; \emph{J. ACM} 48(4), 2001.

\bibitem[DP93]{DP93} C.~Delorme, S.~Poljak. Laplacian eigenvalues and the maximum cut problem.
\emph{Math. Programming} 62, 1993.

\bibitem[PR95]{PR95} S.~Poljak, F.~Rendl. Nonpolyhedral relaxations of graph-bisection
problems. \emph{SIAM J. Optimization} 5, 1995.

\bibitem[Nes98]{Nes98} Y.~Nesterov. Semidefinite relaxation and nonconvex quadratic
optimization. \emph{Optimization Methods and Software} 9, 1998.

\bibitem[AN06]{AN06} N.~Alon, A.~Naor. Approximating the cut-norm via Grothendieck's
inequality. \emph{SIAM J. Computing} 35(4):787--803, 2006.

\bibitem[CW04]{CW04} M.~Charikar, A.~Wirth. Maximizing quadratic programs: extending
Grothen\-dieck's inequality. \emph{Proc.\ 45th IEEE Symposium on Foundations of Computer
Science (FOCS 2004)}, 54--60.

\bibitem[Ban25]{Ban25} A.~Banerjee. On Query Lower Bounds for Aaronson--Kuperberg
Unitary Synthesis Circuits. Research report, 23 July 2025.
\url{https://www.scottaaronson.com/showcase5/synthesis.pdf}.

\bibitem[Hua25]{Hua25} E.~Huang. A 1.5-Query Lower Bound for the Unitary Synthesis Problem.
arXiv:2508.13215v2, 2025.

\bibitem[Bir46]{Bir46} G.~Birkhoff. Tres observaciones sobre el algebra lineal. \emph{Univ.
Nac. Tucum\'an Rev. Ser. A} 5, 1946.

\end{thebibliography}

\end{document}
